% [原文 p.277]
\begin{proof}[证明（续）]
由命题 7.39，得到下述满射：
\[
p_*\cO_S(-K_S-D-\lfloor 2\Delta\rfloor)\to p_*\cO_D(-K_D-\lfloor 2\Delta|_D\rfloor).
\]
由于 $D$ 在 $T$ 上有限，若把 $T$ 收缩到 $p(C)$ 的周围，则 $D$ 上的任意线丛都由全局截面生成。因此可以选取一般元 $D^\sharp_S\in|-K_S-D-\lfloor 2\Delta\rfloor|$，使 $D^\sharp_S\cap C\cap D=\emptyset$。也就是说，$D^\sharp_S$ 在 $C$ 的附近不与 $D$ 相交。因此，在上述任意 $U\ni t$ 上，
\[
D\cdot(D^\sharp_S+\lfloor 2\Delta\rfloor)=D\cdot\lfloor 2\Delta\rfloor\le1
\]
成立。这是因为 $D$ 与 $\Delta$ 只在 $x_t$ 处相交，$D+\Delta$ 是 snc 的，且 $\Delta$ 的任意系数 $a_i$ 满足 $a_i<1$，即 $\lfloor 2a_i\rfloor\le1$。

由构造，$0\sim K_S+D+D^\sharp_S+\lfloor 2\Delta\rfloor$，且由命题 5.58，$(S,D+\lfloor 2\Delta\rfloor+D^\sharp_S)$ 在 $C$ 的附近是 lc 的。
\end{proof}

\begin{propositionp}{7.39}
设 $q\colon Y\to Z$ 是奇点消解。设 $S\subset Y$ 是不可约除子，$B$ 是（不必有效的）Q-除子，且假设 $S+B$ 是 snc 除子。假设 $-(K_Y+S+B)$ 是 $q$-nef 的。此时，下述自然映射是满射：
\[
q_*(\cO_Y(-K_Y-S-\lfloor 2B\rfloor))\to q_*(\cO_S(-K_S-\lfloor 2B|_S\rfloor)).
\]
\end{propositionp}

\begin{proof}
注意下述数值等价关系：
\[
-K_Y-2S-\lfloor 2B\rfloor\equiv K_Y-2(K_Y+S+B)+\{2B\}.
\]
由此，用推论 2.68 得到 $R^1q_*(\cO_Y(-K_Y-2S-\lfloor 2B\rfloor))=0$。考虑下述正合列经 $q$ 的顺像，命题的证明便完成：
\[
0\to\cO_Y(-K_Y-2S-\lfloor 2B\rfloor)\to\cO_Y(-K_Y-S-\lfloor 2B\rfloor)\to\cO_S(-K_S-\lfloor 2B|_S\rfloor)\to0.
\]
\end{proof}

\section{半稳定翻转}
\takeaway{本节给出半稳定翻转的存在定理：在 $s\circ f$ 为 lc、$B$ 有效的假设下，任何关于 $K_X+B$ 的翻转收缩都有翻转；证明分「减去 H」与另一主要阶段，并把 MMP 化为特殊半稳定翻转与除子收缩。}

下面叙述的是本节的主要结果。

\begin{theorem*}{7.40}
设 $(X,B)$ 是三维 klt 对，$B$ 是有效的。设 $f\colon X\to Y$ 是关于 $K_X+B$ 的翻转收缩。假设存在到非奇异曲线 $C$ 的平坦态射 $s\colon Y\to C$，使 $s\circ f$ 是 lc 的。此时，$f$ 有翻转。
\end{theorem*}

\begin{remark}{7.41}
本节的论证在射影的、非固有代数的以及解析的设定下都同样奏效。不过，射影的情形如注 7.55 所述需要稍作修改。
\end{remark}

% [原文 p.278]
\begin{corollaryp}{7.42}
（半稳定翻转的存在）设 $(X,B)$ 是三维 Q-分解 dlt 对，$B$ 是有效的。设 $f\colon X\to Y$ 是关于 $K_X+B$ 的射线型翻转收缩。假设存在到非奇异曲线 $C$ 的平坦态射 $s\colon Y\to C$，使 $s\circ f$ 是 lc 的。此时，$f$ 有翻转。
\end{corollaryp}

\begin{proof}
即使把 $B$ 替换为 $(1-\eps)B$（$0<\eps\ll1$），翻转也是同一个（定义 6.5）。由此，系归约到定理 7.40。
\end{proof}

\par\medskip\noindent\textbf{7.43}\quad 定理 7.40 的证明由定理 7.44 与定理 7.51 两个主要阶段，以及把它们归约下去的容易阶段 7.54 组成。在主要阶段中，作巧妙的状况设定，使那里的 MMP 只靠特殊半稳定翻转与除子收缩就能执行完毕。这一想法最初由 [Sho92] 使用。

下面叙述的是最初的主要阶段，称为「减去 $H$」（subtracting $H$）。\hfill$\square$

\begin{theoremp}{7.44}
设 $(X,B+H)$ 是三维 Q-分解对，$B,H$ 是有效 Q-除子，且 $f\colon X\to Y$ 是射影双有理态射。设 $T\subset Y$ 是有效 Cartier 除子，$S=f^*T$，$S_i$ 是 $S$ 的不可约分支。假设：
\begin{enumerate}[label=(\arabic*),leftmargin=2.4em,itemsep=1pt,topsep=2pt,parsep=0pt]
\item $(X,S+B+H)$ 是 dlt 的。
\item $(X\setminus S,0)$ 是终端的。
\item 存在 $a>0$ 与 $b_j\in\Q$，使
\[
H\equiv_f a(K_X+B)+\sum b_jS_j
\]
成立。
\item $K_X+B+H$ 是 $f$-nef 的。
\end{enumerate}
此时，$(X,B)$ 在 $Y$ 上有极小模型。
\end{theoremp}

\begin{proof}
证明中，在证明若干引理的同时，运行由 $H$ 导出的 $Y$ 上的 MMP。定理 7.44 的记号与假设在这些引理中照原样使用。

% [原文 p.279]
\begin{lemma}{7.45}
（定理 3.5、补足 3.6）存在满足以下 2 个条件的有理数 $\lambda\in[0,1]$。
（1） $K_X+B+\lambda H$ 是 $f$-nef 的。
（2） 若 $\lambda>0$，则存在 $Y$ 上的 $(K_X+B)$-负极射线 $R$，使 $R\cdot(K_X+B+\lambda H)=0$。
\end{lemma}
% [Tutor Note: 题头所引「定理 3.5、补足 3.6」明显错位：本书定理 3.5 是锥定理，第 3 章并无「补足 3.6」。英文原著同样如此，应理解为原书引用编号有误，译文照录。]

若 $\lambda=0$，则定理 7.44 的证明已经完成。假设 $\lambda>0$，并设 $\phi\colon X\to V$ 是 $R$ 的收缩。

\begin{lemmap}{7.46}
$K_X+B$ 在 $Y\setminus T$ 上是 $f$-nef 的，且 $\phi$ 定义了同构态射 $X\setminus S\simeq V\setminus\phi(S)$。
\end{lemmap}

\begin{proof}
前半的主张由条件 (3)、(4) 导出，后半的主张是前半的归结。
\end{proof}

\begin{lemmap}{7.47}
若 $\phi$ 收缩除子 $E$，则上面的条件 (1)–(4) 在把 $f\colon X\to Y$ 换成 $V\to Y$、把 $B,S,H$ 换成 $\phi_*B,\phi_*S,\lambda\phi_*H$ 之后仍然成立。
\end{lemmap}

\begin{proof}
引理由除子收缩的一般性质得到。
\end{proof}

\begin{lemmap}{7.48}
若 $\phi$ 是翻转收缩，则 $\phi$ 是特殊的。设 $p\colon X\dashrightarrow X^+$ 是 $\phi$ 的翻转（由定理 7.32 给出）。此时，上面的条件 (1)–(4) 在把 $f\colon X\to Y$ 换成 $X^+\to Y$、把 $B,S,H$ 换成 $p_*B,p_*S,\lambda p_*H$ 之后仍然成立。
\end{lemmap}

\begin{proof}
必须证明 $\phi$ 是特殊的。首先注意由条件 (2)，$(X\setminus S,0)$ 是终端的。由 $R\cdot(K_X+B+\lambda H)=0$ 与 $R\cdot(K_X+B)<0$ 两个假设，可知 $R\cdot H>0$。由条件 (3)，存在满足 $R\cdot S_{j'}\ne0$ 的 $j'$。进一步，由 $R\cdot S=R\cdot f^*T=0$，若适当地选取 $S_{j''}$，则 $R\cdot S_{j''}$ 与 $R\cdot S_{j'}$ 具有相反符号。这就是 $\phi$ 是特殊的。
\end{proof}

\begin{lemmap}{7.49}
若 $H\ne0$，则无论引理 7.47 与引理 7.48 的哪种情形，都能把上述手续应用于新的设定。如此反复有限次，$H$ 变为 0，并得到 $(X,B)$ 在 $Y$ 上的极小模型。特别地，定理 7.44 成立。
\end{lemmap}

\begin{proof}
虽是巧妙的手续，但注意这是在运行 $(K_X+B)$-MMP。显然除子收缩（引理 7.47）不会无限次发生。引理

% [原文 p.280]
7.48 中的翻转是 $(K_X+B)$-翻转，且前翻转曲线位于 $S$ 上。因此，这不会无限次反复（定理 7.11）。
\end{proof}

这样定理 7.44 的证明就结束了。
\end{proof}

\par\medskip\noindent\textbf{7.50}\quad 第二个主要阶段的主张是：若存在「好的」奇点消解，则在追加修正边界的 MMP 中，可以做到只发生特殊半稳定翻转。经过这一操作，通过把追加的边界减去（定理 7.44），就能到达原 MMP 的目标。半稳定的情形（定理 7.40）中，「好的」奇点消解在 7.54 中构成。\hfill$\square$

% [Tutor Note: 7.50 是全章策略的分水岭：先把边界加到「好消解」，再在追加边界的 MMP 里把边界减回去；额外成分正是定理 7.44 能够适用的原因。]

\begin{theorem*}{7.51}
设 $(X,B)$ 是 $B$ 有效的三维 klt 对，$f\colon X\to Y$ 是到正规 $Y$ 的射影双有理态射，且 $s\colon Y\to C$ 是使 $s\circ f$ 为 lc 的平坦态射。假设：
\begin{enumerate}[label=(\arabic*),leftmargin=2.4em,itemsep=1pt,topsep=2pt,parsep=0pt]
\item $(s\circ f)(\operatorname{Ex}(f))$ 是 0 维的。
\item 存在射影奇点消解 $g\colon Z\to X$ 与 $Z$ 上的自由线性系 $|L|$，满足以下条件：
\begin{enumerate}[label=(\alph*),leftmargin=2.6em,itemsep=1pt,topsep=2pt,parsep=0pt]
\item $s\circ f\circ g\colon Z\to C$ 的每个纤维 $F$ 是约化的，且 $F\cup g_*^{-1}B\cup\operatorname{Ex}(g)$ 的除子部分是 snc 的。
\item 存在 $a>0$、$b_i\in\Q$ 以及含于 $Z\to C$ 的纤维中的不可约除子 $G_i$，使 $L\equiv_{f\circ g}ag^*(K_X+B)+\sum b_iG_i$ 成立。
\end{enumerate}
\end{enumerate}
此时，$(X,B)$ 在 $Y$ 上有标准模型。
\end{theorem*}

\begin{remark}{7.52}
证明中会加入许多约化的边界成分。读者也许觉得奇怪，但请指出：正因如此，才能避免非特殊的半稳定翻转。另外，把定理 7.32(2) 的条件改为关于 $X\setminus(S\cup\operatorname{Supp}\lfloor B\rfloor)$ 而非 $X\setminus S$ 的条件，也是这一手法的影响。
\end{remark}

\begin{proof}
设 $\Sigma\supset(s\circ f)(\operatorname{Ex}(f))\cup g(\bigcup_i G_i)$ 是使 $Z\to C$ 在 $C\setminus\Sigma$ 上光滑的 $C$ 的任意约化除子。记 $T=s^*(\Sigma)$，$S=f^*T$，$G=g^*S$。此时，$S,T,G$ 是约化的（定理 7.51(2.a)），且定理 7.51(2.b) 中出现的 $G_i$ 是 $G$ 的不可约分支。

用没有公共成分的有效 Q-除子 $P,N$，把 $g^*(K_X+B)$ 表示为 $K_Z+P-N$。
% [原文 p.281]
由于 $s\circ f\colon(X,B)\to C$ 是 lc 的，$s\circ f\circ g\colon(Z,P-N)\to C$ 也是 lc 的。因此，$P,G$ 没有公共成分（引理 7.2(2)）。这里注意：$N$ 是 $g$-例外的，且 $N\sim_g K_Z+P$。设 $M\subset Z$ 是满足 $\operatorname{Supp}M=\operatorname{Supp}N\setminus\operatorname{Supp}G$ 的约化除子。此时，对所有的 $\xi\in C\setminus\Sigma$，$(Z,G+P+M+(s\circ f\circ g)^*(\xi))$ 是 dlt 的（定理 7.51(2.a)）。取 $|L|$ 的一般元 $L_1,\cdots,L_4$，记 $H=\sum L_i$，则 $(Z,G+H+P+M)$ 是 dlt 的。

从这样的 $Z$ 出发，运行 $Y$ 上的 $(K_Z+H+P+M)$-MMP。取其中途过程。设某时刻 $Z$ 的像为 $\bar Z$，把 $B,H$ 等在 $\bar Z$ 上的双有理变换记为 $\bar B,\bar H$ 等。记 $h\colon\bar Z\to Y$。
\[
\begin{tikzcd}[column sep=1.3em,row sep=0.5em]
X\setminus S \arrow[d,equal] & \subset & X \arrow[d,"f"] & \bar Z \arrow[l,dashed] \arrow[r,"\phi"] \arrow[dl,"h"'] & W\\
Y\setminus T & \subset & Y \arrow[rr,"s"'] & & C
\end{tikzcd}
\]

\par\medskip\noindent\textbf{主张 7.53}\quad MMP 期间，下述成立。
\begin{enumerate}[label=(\arabic*),leftmargin=2.4em,itemsep=1pt,topsep=2pt,parsep=0pt]
\item $a(K_{\bar Z}+\bar P+\bar M)\equiv_h\bar H+a(\bar M+\bar N)+\sum b_i\bar G_i$（$a>0$，$b_i\in\Q$）。
\item $\bar H$ 是自由的，且 $\bar G\sim_h0$。
\item 设 $\bar F$ 是 $s\circ h$ 的任意纤维。此时，$(\bar Z,\bar G+\bar H+\bar P+\bar M)$、$(\bar Z,\bar F+\bar M)$ 是 dlt 的。
\item 任意射线型收缩 $\phi\colon\bar Z\to W$，是 $\operatorname{Ex}(\phi)\subset\bar G$ 的特殊翻转收缩，或是 $\operatorname{Ex}(\phi)\subset\bar M\cup\bar G$ 的除子收缩。
\item $(\bar Z\setminus(\bar G+\bar M),0)$ 是终端的。
\end{enumerate}
进一步，MMP 期间，翻转存在且终止。因此，MMP 结束，且 $\bar M=0$ 成立。特别地，$\operatorname{Supp}\bar N\subset\operatorname{Supp}\bar G$。

\begin{proof}
(1) 由 $K_Z+P=g^*(K_X+B)+N$ 与定理 7.51(2.b) 得到。

由于 $(\bar Z,\bar H)$ 保持 dlt，$\bar L_1\cap\cdots\cap\bar L_4=\emptyset$（命题 5.63）。因此，$|\bar L|$ 保持自由，且 $\bar H$ 也是自由的。又 $\bar G=h^*(T)$。这就得到 (2)。这里运行的既是 $(K_Z+G+H+P+M)$-MMP（注意 $G\sim_h0$），同时也是 $(K_Z+F+P+M)$-MMP（注意 $\bar H$ 是 nef 的）。由此得到 (3)。

% [原文 p.282]
为证 (4)，设 $\ell\subset\operatorname{Ex}(\phi)$ 是任意不可约曲线。由 (1) 与 (2) 可知 $a\ell\cdot(\bar M+\bar N)+\sum b_i(\ell\cdot\bar G_i)<0$ 成立。

首先，对所有 $\bar G_i$ 设 $\ell\cdot\bar G_i=0$。此时，对适当的不可约分支 $\bar M_k\subset\bar M$ 有 $\ell\cdot\bar M_k<0$。特别地，$\ell\subset\bar M_k$。进一步，$\bar M_k$ 是正规的（推论 5.52）。设 $\bar F$ 是包含 $\ell$ 的 $s\circ h$ 的纤维，$\bar F_j\supset\ell$ 是 $\bar F$ 的不可约分支。此时，得到 $(\ell\cdot\bar F_j)=0$，这是因为当 $\bar F\not\subset\bar G$ 时 $\bar F_j=\bar F$，而当 $\bar F\subset\bar G$ 时由 $\ell\cdot\bar G_i=0$ 导出。由于 $(\bar Z,\bar F+\bar M)$ 是 dlt 的，$\bar F_j$ 在余维 2 处是 Cartier 的（推论 5.56）。进一步，因为 $\ell\subset\bar F_j\cap\bar M_k$，作为 $\ell$ 附近的集合有 $\bar F_j|_{\bar M_k}=\ell$（推论 5.62）。因此，$\ell$ 是 $\bar M_k$ 上的 Q-Cartier 除子，且满足 $(\ell^2)_{\bar M_k}=0$。由此，$\bar M_k\to\phi(\bar M_k)$ 不是双有理的（引理 3.39），故得到 $\operatorname{Ex}(\phi)=\bar M_k$。

其次，假设存在满足 $\ell\cdot\bar G_i\ne0$ 的 $\bar G_i$。由于 $\bar G=\sum\bar G_i\sim_h0$，还存在另一个 $\bar G_j$ 使 $(\ell\cdot\bar G_i)(\ell\cdot\bar G_j)<0$。因此，$\operatorname{Ex}(\phi)\subset\bar G_i\cup\bar G_j\subset\bar G$。若 $\phi$ 是翻转收缩，则是特殊的。这就证明了 (4)。

由 (4)，$\bar Z\setminus(\bar G+\bar M)$ 在运行 MMP 期间保持同型。由于 $Z$ 非奇异，故导出 (5)。特殊半稳定翻转存在（定理 7.32），而且终止（定理 7.7）。MMP 结束时，$K_{\bar Z}+\bar H+\bar P+\bar M$ 是 $h$-nef 的。由定理 7.51(2.b) 与 $X\setminus S\simeq Y\setminus T$，注意 $\bar H|_{\bar Z\setminus\bar G}\equiv_h0$，便知 $(\bar M+\bar N)|_{\bar Z\setminus\bar G}\equiv_h(K_{\bar Z}+\bar P+\bar M)|_{\bar Z\setminus\bar G}$ 在 $Y\setminus T$ 上是 $h$-nef 的。由于它是 $\bar Z\setminus\bar G$ 上的有效 $h$-例外除子，故可知 $\bar M=0$（引理 3.39）。
\end{proof}

设 $p\colon W\to Y$ 是极小模型。用 $G_W,H_W$ 等表示 $G,H$ 等的双有理变换。由主张 7.53 可知下述事实。$G_W=p^*T$，$\operatorname{Supp}N_W\subset\operatorname{Supp}G_W$，$(W\setminus G_W,0)$ 是终端的，$(W,G_W+P_W+H_W)$ 是 dlt 的，$K_W+P_W+H_W$ 是 $p$-nef 的，并且 $H_W$ 是自由的。这些满足定理 7.44 的条件。

因此，由定理 7.44，$(W,P_W)$ 在 $Y$ 上有极小模型。由于 $(X,B)$ 是 klt 的，得到 $\lfloor P\rfloor=0$，特别地 $\lfloor P_W\rfloor=0$。因此，把相对无基点定理 3.24 应用于极小模型，便得到 $(W,P_W)$ 的标准模型。
% [原文 p.283]
由主张 7.53(2)，注意 $(K_Z+H+P+M)$-MMP 也是 $(K_Z+P+M)$-MMP 的一部分。因此特别地，$(W,P_W)$ 与 $(Z,P+M)$ 在 $Y$ 上有同一标准模型。在 $X$ 的奇点消解 $Z$ 上，$K_Z+P+M=g^*(K_X+B)+M+N$ 且 $N\subset\operatorname{Ex}(g)$ 成立，所以它也是 $(X,B)$ 在 $Y$ 上的标准模型（推论 3.53）。
% [Tutor Note: 此处引「推论 3.53」与用词不谐：该推论讲的是伪标准模型的唯一性，而结论是标准模型的存在。原文作「推论 3.53」，疑为原书引用疏失，译文照录并以「存在—唯一」分别对应定理 7.51 与推论 3.53 理解。]

% [Tutor Note: 定理 7.51 的证明分两步：先用半稳定约化造出好消解并把边界补成 snc，再对修正后的边界运行 MMP；主张 7.53 保证全程只出现特殊半稳定翻转。]
\end{proof}

\par\medskip\noindent\textbf{7.54}（定理 7.40 的证明）\quad 设 $a>0$ 是使 $Q=a(K_X+B)$ 成为 Cartier 的整数。把 $Y$ 收缩为 $f(\operatorname{Ex}(f))$ 的邻域，使 $f_*\cO_X(Q)$ 由全局截面生成。设 $Bs|Q|$ 是 $|Q|$ 的所有元作为概形的公共部分。此时，$\operatorname{Supp}Bs|Q|\subset\operatorname{Ex}(f)$，且 $\operatorname{Supp}Bs|Q|$ 的余维为 2 以上。

$g_1\colon X_1\to X$ 是 $X$ 沿 $Bs|Q|$ 的爆破的正规化。存在 Cartier 除子 $E_1\subset\operatorname{Supp}(g_1)^{-1}\operatorname{Ex}(f)$ 使 $|L_1|=g_1^*|Q|-E_1$ 自由。

选取充分分歧的有限覆盖 $r\colon C'\to C$，把半稳定约化定理 7.17、即定理 7.19 应用于 $X_1\to C$ 与部分集合 $\operatorname{Ex}(g_1)\cup(g_1)^{-1}(B)$。设 $Y',X',X_1'$ 分别是 $Y,X,X_1$ 经 $r$ 的基变换（的正规化）。设 $Z'$ 是定理 7.17 给出的 $X_1'$ 的奇点消解。于是得到由诱导态射组成的如下交换图。
\[
\begin{tikzcd}[column sep=1.1em,row sep=1.5em]
Z' \arrow[r,"g_2'"] \arrow[dr,"p_1"'] & X_1' \arrow[r,"g_1'"] \arrow[d,"p"'] & X' \arrow[r,"f'"] \arrow[d,"q"'] & Y' \arrow[r,"s'"] \arrow[d,"r"'] & C'\\
& X_1 \arrow[r,"g_1"'] & X \arrow[r,"f"'] & Y \arrow[r,"s"'] & C
\end{tikzcd}
\]
然后，记 $g'\colon Z'\to X'$、$B'=p^*B$、$|L'|=(p_1\circ g_2')^*|L_1|$。

其次，验证 $(X',B')$、$f'\colon X'\to Y'$、$|L'|$ 满足定理 7.51 的条件。由于 $(X,B)$ 是 klt 且 $(s\circ f)\colon(X,B)\to C$ 是 lc，$(s'\circ f')\colon(X',B')\to C'$ 是 lc（引理 7.6），且 $(X',B')$ 是 klt（引理 7.2）。由于 $\operatorname{Ex}(f')\subset p^{-1}(\operatorname{Ex}(f))$，定理 7.51(1) 得到满足。定理 7.51(2.a) 由定理 7.17 导出。$|L'|$ 的自由性由 $|L_1|$ 的自由性可知。由构造得到次数关系
\[
L'\sim a(p\circ g')^*(K_X+B)-(p_1\circ g_2')^*E_1.
\]

% [原文 p.284]
注意 $(p_1\circ g_2')^*E_1$ 与 $(g')^*(K_{X'}+B')-(p\circ g')^*(K_X+B)$ 是 $Z'\to C'$ 的纤维的不可约成分的线性组合的线性等价（命题 7.23(3)）。这样就得到定理 7.51(2.b)，从而可以应用定理 7.51。

于是，$(X',B')$ 在 $Y'$ 上有标准模型，它是 $f'$ 的翻转（推论 6.4），且 $f$ 有翻转（推论 6.9）。\hfill$\blacksquare$

\begin{remark}{7.55}
若 $X,Y,C$ 全是射影簇，且只证明关于射影簇的论断，则 7.54 中只有把簇收缩到邻域的讨论部分需要修正，本节其余部分不需要修正。

按如下做即可。在 7.54 中定义 $Q$ 之前，选取 $Y$ 上的丰富除子 $A$，使 $f_*\cO_X(a(K_X+B))\otimes\cO_Y(A)$ 由全局截面生成，不收缩 $Y$，定义 $Q=a(K_X+B)+f^*A$。其余部分无需修正即可照原样起作用。
\end{remark}

\section{对曲面族的应用}
\takeaway{本节把半稳定 MMP 用于三个问题：一般型曲面模空间的紧化、曲面奇点半普遍形变空间的研究，以及曲面族的同时奇点消解；核心结果是给出同时奇点消解存在的充要条件。}

本节的目的是讨论半稳定 MMP 的三个应用。首先讨论一般型曲面模空间的紧化。回顾历史，成为半稳定 MMP 研究出发点的正是与此相关的问题。

其次，讨论曲面奇点 $H$ 的半普遍形变空间 $\operatorname{Def}(H)$ 的研究。在商奇点的情况下，可以用半稳定 MMP 研究 $\operatorname{Def}(H)$ 的复杂结构。

最后，研究曲面族的同时奇点消解。作为主要结果，给出同时奇点消解存在的充要条件。

\par\medskip\noindent\textbf{7.56}（模空间的紧化）\quad 长期以来已知：由所有非奇异射影簇的给定数值不变量构成的集合，与某簇（所谓模空间）的点全体之间，存在自然意义下的一一对应。最简单的例子是椭圆曲线。此时，$j$-不变量在 $\C$ 上的椭圆曲线与 $\C$ 的点之间建立一一对应。给定亏格 $g$ 的曲线对应到由通常 $M_g$ 表示的簇的点，这也是已知的。
% [原文 p.285]
由于空间 $M_g$ 不是固有的，寻找它的自然紧化便作为饶有兴味的问题而受到关注。由此产生下述问题：「与 $M_g$ 的边界对应的，是什么样的代数几何对象？」

一个答案由 [DM69] 确立。

\begin{definition}{7.57}
域 $k$ 上亏格 $g$ 的\textbf{稳定曲线}（stable curve）是指满足下述条件的、域 $k$ 上的 1 维固有概形 $D$：奇点至多是通常二重点，$\omega_D$ 丰富，且 $\chi(D,\cO_D)=1-g$。

\textbf{稳定曲线族}是指各纤维都是稳定曲线的固有且平坦的态射 $f\colon X\to S$。
\end{definition}

% [Tutor Note: 稳定曲线把「边界」限定为通常二重点，因而模空间 $\overline{M}_g$ 紧；书中随后给出允许结点与尖点的替代定义，两条路线给出不同但都良好的紧化，读时注意区分。]

亏格 $g$ 的任意稳定曲线，与把 $M_g$ 作为稠密开子簇包含的 $\overline{M}_g$ 的任意点自然一一对应 [DM69]。

应当强调，紧化问题还有别的解答。例如，考虑曲线 $D$：奇点仅为通常结点与通常尖点，$\omega_D$ 丰富，$\chi(D,\cO_D)=1-g$。不过需要追加条件：不存在这样的 $D$ 的亏格 1 的不可约分支，它与其余所有分支仅在一点相交。即使从这个定义出发，也能建立良好的理论。

还有对高维簇建立类似理论的问题。根据经验，作为 7.56 特殊情形的下述问题最为重要。

设 $C$ 是非奇异曲线，$C^0\subset C$ 是开子集。假设给定曲线族 $f^0\colon X^0\to C^0$，其所有纤维都是亏格 $g$ 的非奇异曲线。如何把它延拓为族 $f\colon X\to C$？

某种意义上，若对 $C$ 为离散赋值环 $V$ 的谱 $\operatorname{Spec}V$ 的情形找到满意的解答，则 7.56 的一般情形也能解决。（试与用分离性与固有性的赋值判别法比较。）\hfill$\square$

下述引理说明这个问题如何与 MMP 相关联。（证明简单，交给读者。参看引理 7.59。）

\begin{lemma*}{7.58}
设从曲面到非奇异曲线的态射 $g\colon S\to C$ 的一般纤维非奇异。此时下述等价。
\begin{enumerate}[label=(\arabic*),leftmargin=2.4em,itemsep=1pt,topsep=2pt,parsep=0pt]
\item 任何纤维至多只有通常结点。
\item 若 $C'$ 是非奇异曲线，$h\colon C'\to C$ 是非常数态射，则 $S\times_C C'$ 是标准的。
\item $g\colon S\to C$ 是 lc 的。
\end{enumerate}
\end{lemma*}

% [原文 p.286]
条件 7.58(2)–(3) 在所有维数下都有意义，且等价。

\begin{lemmap}{7.59}
设 $g\colon X\to C$ 是从簇到非奇异曲线的态射，一般纤维是标准的。此时下述等价。
\begin{enumerate}[label=(\arabic*),leftmargin=2.4em,itemsep=1pt,topsep=2pt,parsep=0pt]
\item 若 $C'$ 是非奇异曲线，$h\colon C'\to C$ 是非常数态射，则 $X\times_C C'$ 是标准的。
\item $g\colon X\to C$ 是 lc 的。
\end{enumerate}
\end{lemmap}

\begin{proof}
(2)$\Rightarrow$(1) 由引理 7.2 得到。为证逆，取点 $c\in C$，选择态射 $h\colon C'\to C$ 使 $c'=h^{-1}(c)$ 为一点、且是 $h$ 的唯一分歧点，记其分歧指数为 $r$。此时，
\[
K_{X'}=h^*\left(K_X+\left(1-\frac1r\right)g^{-1}(c)\right).
\]
由假设 $X'$ 是标准的，特别地是 klt 的。因此由命题 5.20，$(X,(1-\frac1r)g^{-1}(c))$ 是 klt 的。由于 $r>0$ 可以任意大，由推论 2.35 得 $(X,g^{-1}(c))$ 是 lc 的。
\end{proof}

% [Tutor Note: 证明的要点是用充分分歧的覆盖把纤维条件提升为整体的 lc 条件，再令分歧指数趋于无穷；这是全书中反复出现的「基变换检验」手法。]

\par\medskip\noindent\textbf{7.60}\quad 引理 7.58(1) 在高维中的类似问题仍然存在。伴随反转（定理 5.50）所暗示的正确条件，至少对正规纤维而言是「各纤维都是 lc 的」。

非正规的情形需要 lc 概念的适当推广。这一概念称为半对数标准的（semi-log canonical），或简称 slc。一旦基本问题得到解决，其理论与 lc 情形非常相似。详见 [KSB88, K+92]。

然而，引理 7.58(1) 在高维中没有类似物。态射是否为 lc 不能仅由纤维刻画。这从下述例子可以看出。

\par\medskip\noindent\textbf{例 7.61}\quad 从簇到超平面截面上的锥的退化，按如下方式构成。设 $Z\subset\PP^n$ 是射影正规（projectively normal）簇，$D=(x_n=0)$ 是 $Z$ 的超平面截面。
